\subsection{Горизонтальная сила трения}
%%%%%%%%%%%%%%%%%%%%%%%%%%%%%%%%%%%%%%%%%%%%%%%%%%%%
%%%%%%%%%%%%%%%%%%%%%%%%%%%%%%%%%%%%%%%%%%%%%%%%%%%% 
 \z[{\begin{wrapfigure}[5]{r}{0.3\textwidth}%
		\vspace{-15pt}%
		\begin{flushright}%
 			\includegraphics[width=0.3\textwidth]{5.2.1}%
		\end{flushright}%
	\end{wrapfigure}%
	}]%
	{В системе, изображённой на рисунке, тело массой $M$ может скользить без трения по горизонтальной плоскости. Коэффициент трения между телами $M$ и $m$ равен $\mu$. Найдите ускорение $a$ тела $M$. Массой блоков и нерастяжимой нити пренебречь. Ускорение свободного падения равно $g$.}
 %
{$a=\frac{2mg}{M+(5+2\mu)m}$}
%%%%%%%%%%%%%%%%%%%%%%%%%%%%%%%%%%%%%%%%%%%%%%%%%%%%	
 \z[{\begin{wrapfigure}[4]{r}{0.3\textwidth}%
		\vspace{-15pt}%
		\begin{flushright}%
 			\includegraphics[width=0.3\textwidth]{5.2.2}%
		\end{flushright}%
	\end{wrapfigure}%
	}]%
	{На горизонтальной гладкой поверхности стола покоится доска массой $M$ (см. рисунок). На доску со скоростью $v$ въезжает шайба массы $m$. Какой должна быть длина доски, чтобы шайба не соскользнула с неё? Коэффициент трения скольжения между шайбой и доской равен $\mu$, размер шайбы мал по сравнению с длиной доски.}
 %
{$L>\frac{mv^2}{2\mu(M+m)g}$}
%%%%%%%%%%%%%%%%%%%%%%%%%%%%%%%%%%%%%%%%%%%%%%%%%%%%	
 \z[{\begin{wrapfigure}[5]{r}{0.3\textwidth}%
		\vspace{-12pt}%
		\begin{flushright}%
 			\includegraphics[width=0.3\textwidth]{5.2.3}%
		\end{flushright}%
	\end{wrapfigure}%
	}]%
	{Доска массы $M$ и длины $L$ скользит с некоторой скоростью $v_0$ по гладкой горизонтальной поверхности. На левом краю доски лежит кубик массы $m$. Коэффициент трения скольжения между кубиком и доской равен $\mu$. Доска испытывает абсолютно упругий удар о вертикальную стенку (см. рисунок). При какой максимальной скорости доски $v_0  = v_{\max}$ кубик с неё не упадёт? Размерами кубика по сравнению с $L$ пренебречь. В процессе всего движения кубик не опрокидывается.}
 %
{$v_{\max}=\sqrt{\frac{1}{2}\mu g L \left(1+\frac{m}{M}\right)}$}
%%%%%%%%%%%%%%%%%%%%%%%%%%%%%%%%%%%%%%%%%%%%%%%%%%%%
 \z{Человек массы $m_1$, оставаясь на месте, тянет за веревку груз массы $m_2$. Коэффициент трения о горизонтальную плоскость равен $\mu$. При какой наименьшей силе натяжения веревки груз стронется с места? Под каким углом к горизонтальной плоскости должна быть направлена веревка?}
%
{$T_{\min}=\frac{\mu m_2g}{\sqrt{1+\mu^2}};~\tg \alpha = \mu$ при $m_1>m_2 \frac{1-\mu^2}{1+\mu^2}$ \par
$T_{\min} = \mu g \sqrt{\frac{m_1^2+m_2^2}{2}};~ \tg \alpha = \frac{m_2-m_1}{\mu (m_1+m_2)}$ при $m_1>m_2 \frac{1-\mu^2}{1+\mu^2}$
}
%%%%%%%%%%%%%%%%%%%%%%%%%%%%%%%%%%%%%%%%%%%%%%%%%%%%	
 \z[{\begin{wrapfigure}[4]{r}{0.3\textwidth}%
		\vspace{-10pt}%
		\begin{flushright}%
 			\includegraphics[width=0.3\textwidth]{5.2.5}%
		\end{flushright}%
	\end{wrapfigure}%
	}]%
	{На гладком горизонтальном столе расположена система грузов, изображенная на рисунке. Правый нижний груз тянут вдоль стола с силой $F$, как указано на рисунке. Коэффициент трения между грузами массы $m_1$ и $m_2$ равен $\mu$. Найдите ускорение всех грузов системы.}
 %
{При $F<\frac{2\mu m_1g(m1+m2)}{m_2+2m_1}=F_0$: $a_\text{1лев} =a_\text{1прав} = a_\text{2прав} = \frac{F}{2(m_1+m_2)}$ \par
при $F>F_0$: $a_\text{2прав}=\frac{F-\mu m_1g}{m_2}, ~a_\text{1лев} =a_\text{1прав} = a_\text{2лев} = \frac{\mu m_1g}{2m_1+m_2} $}